\documentclass[11pt]{article}

\usepackage[margin=1in]{geometry}
\usepackage{amsmath}
\usepackage{booktabs}
\usepackage[T1]{fontenc}
\usepackage{lmodern}
\usepackage[hidelinks]{hyperref}

\title{IEEE 9-Bus AC N-1 Branch-Outage Analysis}
\author{Load Flow Lite}
\date{}

\newcommand{\pu}{\mathrm{p.u.}}
\newcommand{\MW}{\mathrm{MW}}
\newcommand{\MVAr}{\mathrm{MVAr}}

\begin{document}

\maketitle

\begin{abstract}
This note documents an AC N-1 branch-outage study for the IEEE 9-bus case
in \path{data/case9.json}. Each branch is removed one at a time, the
remaining AC power-flow problem is classified topologically, and solvable
post-contingency cases are solved with the Newton--Raphson solver in this
repository. The checked-in JSON report is
\path{data/case9_n1_results.json}.
\end{abstract}

\section{Scope}

N-1 analysis asks whether the system remains acceptable after the loss of
one component. Here the components are the nine branch records in the
IEEE 9-bus fixture. The study is deliberately narrow:

\begin{itemize}
  \item AC power flow is solved for connected post-contingency networks and
  for generator-islanding cases whose surviving component still contains
  the original slack bus and all load.
  \item No optimal power flow, economic redispatch, dynamic stability,
  transient stability, protection action, or load shedding is modeled.
  \item Branch thermal ratings and generator reactive limits are absent
  from the fixtures. The report therefore cannot certify thermal-loading
  or reactive-limit compliance.
  \item PV-to-PQ conversion is not performed. Solved generator reactive
  outputs are reported as diagnostics only.
\end{itemize}

The voltage screening band is $0.95 \le |V| \le 1.05~\pu$. All power-flow
solver quantities are on the 100 MVA base unless explicitly shown in MW or
MVAr.

\section{Topology and Contingencies}

The branches split into the six-branch transmission loop and three radial
generator-connection branches shown in Table~\ref{tab:classification}. The
stored branch orientation is used as the branch identifier.

\begin{table}[htbp]
\centering
\caption{Branch classification used for the case9 N-1 study.}
\label{tab:classification}
\begin{tabular}{@{}ll@{}}
\toprule
Group & Branches \\
\midrule
Transmission loop & 4--5, 5--6, 6--7, 7--8, 8--9, 9--4 \\
Generator connection & 1--4, 3--6, 8--2 \\
\bottomrule
\end{tabular}
\end{table}

Removing branch 1--4 isolates the slack bus. The remaining load component
has no angle reference, so the ordinary AC power-flow problem is not
formulated. Removing branch 3--6 or 8--2 isolates a single PV generator
bus. In this implementation the surviving slack-containing component is
solved, with the missing real generation absorbed by the original slack
bus. That is a modeling choice for this small study, not a general
graph-theoretic requirement.

\section{Base Case}

The base case converges in four Newton iterations. Its solved voltage range
is $0.9956~\pu$ at bus 9 to $1.0400~\pu$ at bus 1. The real-power loss is
$4.641~\MW$. The branch reactive endpoint sum,
\[
  \sum_{\ell} (Q_{\ell,\mathrm{from}} + Q_{\ell,\mathrm{to}}),
\]
is $-92.160~\MVAr$. This quantity includes the effect of line charging and
should not be read as colloquial reactive losses.

\section{Solved Contingencies}

Table~\ref{tab:solved} summarizes every solved connected or partial-island
case. The largest branch-end flow column is the largest absolute real-power
flow at either endpoint of any surviving branch.

\begin{table}[htbp]
\centering
\caption{Solved N-1 results.}
\label{tab:solved}
\small
\begin{tabular}{@{}lcrrrrr@{}}
\toprule
Outage & Status & Min $|V|$ & Bus & Loss & $Q_f+Q_t$ & Max $|P|$ \\
 & & p.u. & & MW & MVAr & MW \\
\midrule
4--5 & solved & 0.9418 & 5 & 6.136 & -63.274 & 163.000 \\
5--6 & solved & 0.9639 & 5 & 9.491 & -21.579 & 163.000 \\
3--6 & partial island & 1.0028 & 9 & 3.629 & -88.881 & 163.000 \\
6--7 & solved & 0.9783 & 7 & 5.353 & -62.616 & 163.000 \\
7--8 & solved & 0.9690 & 7 & 12.093 & -28.468 & 163.000 \\
8--2 & partial island & 0.9924 & 9 & 4.435 & -71.541 & 234.435 \\
8--9 & solved & 0.9380 & 9 & 13.208 & 0.759 & 163.000 \\
9--4 & solved & 0.8388 & 9 & 9.567 & -38.782 & 163.000 \\
\bottomrule
\end{tabular}
\end{table}

The largest voltage changes from the base case occur at bus 9 for outage
9--4, where $|V|$ falls by $0.1569~\pu$, and at bus 5 for outage 4--5,
where $|V|$ falls by $0.0708~\pu$. The 3--6 and 8--2 partial-island cases
remain close to the base voltage profile because the transmission loop and
all loads remain connected to the original slack bus.

\section{Islanding and Voltage Violations}

The generator-connection outages are summarized in Table~\ref{tab:islands}.

\begin{table}[htbp]
\centering
\caption{Generator-connection outage treatment.}
\label{tab:islands}
\begin{tabular}{@{}lll@{}}
\toprule
Outage & Status & Treatment \\
\midrule
1--4 & slack islanded & Not solved; load component has no slack bus \\
3--6 & partial island & Bus 3 removed; 85 MW absorbed by slack \\
8--2 & partial island & Bus 2 removed; 163 MW absorbed by slack \\
\bottomrule
\end{tabular}
\end{table}

Voltage-limit violations are low-voltage violations only:

\begin{table}[htbp]
\centering
\caption{Post-contingency voltage violations for the 0.95--1.05 p.u. band.}
\label{tab:violations}
\begin{tabular}{@{}lrr@{}}
\toprule
Outage & Bus & $|V|$ p.u. \\
\midrule
4--5 & 5 & 0.9418 \\
8--9 & 9 & 0.9380 \\
9--4 & 9 & 0.8388 \\
\bottomrule
\end{tabular}
\end{table}

No solved generator reactive output exceeds the configured diagnostic
threshold of $2.0~\pu$ on the system base. This is only a soft diagnostic:
the fixture has no $Q_{\min}$ or $Q_{\max}$ limits, so no reactive-limit
compliance claim is made.

\section{Verification}

The implementation checks the base case against the existing MATPOWER
fixture and checks the six connected transmission-loop outages against the
static MATPOWER fixture in \path{data/case9_n1_solutions.json}. For each
solved or partial-island case, active and reactive balance are checked on
the surviving component:
\[
  \sum_i (P_{G,i} - P_{D,i}) = \sum_{\ell} P_{\ell,\mathrm{loss}},
  \qquad
  \sum_i (Q_{G,i} - Q_{D,i}) =
  \sum_{\ell} (Q_{\ell,\mathrm{from}} + Q_{\ell,\mathrm{to}}).
\]
The JSON output is deterministic by default and omits timestamps unless the
script is called with \texttt{--timestamp}.

\end{document}
